\chapter{Dataset}
\label{ch:dataset}

\section{Synthetic Pharmacotherapy Dataset}

The primary dataset used in this thesis is Synthetic Pharmacotherapy Graph v3, a controlled benchmark designed to connect a known mechanistic graph structure with patient-level empirical observations.
The dataset was constructed specifically to support both tasks investigated in the thesis: reconstruction of graph relations in Task~A and patient-level endpoint prediction in Task~B.
It consists of two complementary components: a directed pharmacotherapy graph and patient-level records generated from that graph.

The use of synthetic data provides access to a known reference graph and to the complete data-generating process.
This makes it possible to introduce controlled observational artefacts while preserving the underlying mechanistic structure.
The dataset should therefore be interpreted as methodological ground truth for model evaluation.
Graph coefficients, baseline probabilities and realised endpoint rates are simulation parameters and must not be interpreted as estimates of real clinical effects.

\section{Pharmacotherapy Knowledge Graph}

The graph was designed as a directed acyclic graph (DAG) representing pathways from patient characteristics and drug exposures, through physiological mechanisms and intermediate adverse states, to clinical outcomes.
Graph construction started from clinically relevant endpoints and proceeded backwards through their expected determinants and mechanistic pathways.

Domain knowledge in pharmacology was used to define and review the initial set of drug classes, risk factors, mechanisms and clinical endpoints.
Claude Fable~5 was used as an AI-assisted modelling tool to propose candidate mechanistic pathways, including drug--drug and drug--disease interaction structures.
Importantly, the model was used to support pathway drafting rather than to determine the final graph automatically.
Proposed relations were reviewed and could be accepted, rejected or modified before inclusion in the graph specification.

The graph was developed iteratively and subsequently audited for structural consistency.
The final v3 graph contains 155 nodes and 405 directed edges and is acyclic, with no duplicate nodes, duplicate edges or references to missing nodes.
The nodes are organised into six semantic categories:

\begin{table}[ht]
\centering
\caption{Node types in Synthetic Pharmacotherapy Graph v3.}
\label{tab:dataset-node-types}
\begin{tabular}{@{}llp{7.2cm}@{}}
\toprule
Node type & Number & Role \\
\midrule
\texttt{patient\_context} & 24 & demographics, comorbidities and other patient-level characteristics \\
\texttt{drug\_exposure} & 29 & exposure to drug classes \\
\texttt{mechanism} & 49 & physiological mechanisms, interaction gates and drug-load variables \\
\texttt{adr\_or\_intermediate\_state} & 28 & intermediate adverse or physiological states \\
\texttt{observation\_or\_selection} & 12 & testing, documentation and selection processes \\
\texttt{clinical\_endpoint} & 13 & target clinical outcomes \\
\bottomrule
\end{tabular}
\end{table}

These categories distinguish patient context, treatment, biological mechanisms, intermediate consequences, observation processes and final clinical events.

\begin{figure}[htbp]
\centering
\includegraphics[width=\textwidth]{01_node_counts_by_type_and_layer.png}
\caption{Number of nodes in each semantic layer of Synthetic Pharmacotherapy Graph~v3.}
\label{fig:dataset-node-counts}
\end{figure}

Directed edges represent different semantic relationships, including background risk, treatment assignment, drug-to-mechanism effects, mechanistic progression, adverse-state progression, observation processes, and drug--drug or drug--disease interactions.
Edge attributes include the simulated magnitude and sign of an effect, which are subsequently used by the patient generator.

The graph contains 13 clinical endpoints regarded as points of clinical state requiring intervention:
AKI, DILI, Depression, Falls, Delirium, GI bleeding, Hyponatremia, Hyperkalemia, QT arrhythmia, Hospitalization, Serotonin syndrome, Rhabdomyolysis, and Lactic acidosis.

\begin{figure}[htbp]
\centering
\includegraphics[width=\textwidth]{causal_dag_type_overview.pdf}
\caption{Overview of the synthetic pharmacotherapy DAG with node colours by semantic type.}
\label{fig:dag-overview}
\end{figure}

\begin{figure}[htbp]
\centering
\includegraphics[width=\textwidth]{fig_aki_kg_fragment_networkx.png}
\caption{AKI-related fragment of the Synthetic Pharmacotherapy Graph~v3: all audited ancestors of the AKI endpoint ($32$ nodes / $79$ edges), drawn with NetworkX.
Node colours follow the v3 ontology (patient context, drug exposure, mechanism, ADR/intermediate state, clinical endpoint).
The full graph is larger; this excerpt is for orientation.}
\label{fig:aki-kg-fragment}
\end{figure}

\section{Patient-Level Data Generation}

The graph defines a common mechanistic structure, whereas the patient table represents individual realisations of that structure.
For each simulated patient, values are generated for the graph variables in topological order.
Patient-context variables are generated first, followed by drug exposures, mechanisms, intermediate states and finally clinical endpoints.
As a result, a parent variable is always realised before a downstream child variable is generated.

Most stochastic binary nodes are generated using an additive logistic model.
For node $i$,
\begin{equation}
\eta_i
=
\mathrm{logit}(p_i^{0})
+
\sum_{j\in\mathrm{parents}(i)}
s_{ij}\,w_{ij}\,x_j,
\end{equation}
followed by
\begin{equation}
x_i\sim\mathrm{Bernoulli}\!\left(\sigma(\eta_i)\right),
\end{equation}
where $p_i^{0}$ denotes the baseline probability of the node, $x_j$ is the realised value of parent $j$, and $w_{ij}$ and $s_{ij}$ denote the simulated effect magnitude and sign associated with the edge $j\to i$.
Count and continuous parents are standardised before entering the additive predictor.

Not all variables are stochastic.
The generator also contains deterministic quantities representing clinically meaningful combinations such as drug loads, drug--drug interaction gates, drug--disease interactions and cumulative adverse-event burden.
For example, some interaction variables are activated only when several required drug exposures are simultaneously present.

The dataset additionally separates an underlying clinical state from whether that state is observed or documented.
For example,
\begin{equation}
\texttt{elevated\_creatinine\_recorded}
=
\texttt{creatinine\_rise}
\times
\texttt{creatinine\_tested}.
\end{equation}
A true creatinine rise can therefore occur without being present in the recorded representation if testing is not performed.
This distinction allows documentation and selection artefacts to be introduced without changing the underlying mechanistic graph.

The clean dataset contains 20{,}000 simulated patients.
Each row corresponds to one patient, while graph nodes correspond to patient-level variables.
The 155 graph variables include binary, count and continuous values; the dataset additionally contains technical identifiers such as \texttt{patient\_id} and \texttt{hospital\_id}.

\section{Observational Scenarios}

Six scenarios were generated from the same underlying pharmacotherapy graph.
They were designed to test whether the models remain robust when the patient-level evidence is affected by different observational artefacts.

\begin{table}[ht]
\centering
\caption{Observational scenarios derived from the same pharmacotherapy graph.}
\label{tab:dataset-scenarios}
\begin{tabularx}{\textwidth}{@{}lXr@{}}
\toprule
Scenario & Purpose & Patients \\
\midrule
\texttt{clean} & reference setting with complete generated information & 20{,}000 \\
\texttt{hidden\_confounder} & latent severity remains active in generation but is withheld from observed data & 20{,}000 \\
\texttt{selection\_bias} & cohort restricted according to healthcare-contact and testing variables & 14{,}062 \\
\texttt{no\_overlap} & selected treatment exposures become rare in particular patient subgroups & 20{,}000 \\
\texttt{noisy\_documentation} & imperfect recording with false-negative and false-positive documentation & 20{,}000 \\
\texttt{multihospital} & heterogeneous treatment and documentation processes across four simulated hospitals & 20{,}000 \\
\bottomrule
\end{tabularx}
\end{table}

The scenarios differ in the observability and distribution of patient-level data, while the underlying graph semantics are preserved.
Importantly, the \texttt{selection\_bias} dataset is an intentional subset of the original population and therefore contains 14{,}062 patients.

\section{Patient Splits and Use Across Tasks}

Patients are divided into training, validation and test sets in proportions of 70\%, 15\% and 15\%.
For the clean cohort, this corresponds to 14{,}000 training patients, 3{,}000 validation patients and 3{,}000 test patients.
The split is defined at the patient level and reused across the alternative scenarios to maintain comparable patient partitions whenever possible.

The two tasks use this split differently.
In Task~A, patient records are not themselves the prediction instances: the training-patient partition is used to estimate node-level statistics and pairwise dependence measures, while the prediction instances are candidate graph edges.
The Task~A candidate inventory, negative-sampling rule and edge-level split are defined in Chapter~\ref{ch:taskA-rozwiazanie}.

In Task~B, the patient split itself defines the training, validation and test populations.
Each patient contributes individual node values and clinical endpoint labels, while the graph provides the common mechanistic scaffold.

All statistics derived from patient-level data are computed using the appropriate training partition only.
For Task~B, endpoint values, post-event or definitional intermediates, and other variables unavailable at the intended prediction point are excluded from predictive inputs.
